% !TeX root = ../main.tex
\section{Evaluation of the Proposed Offloading Methods}
\label{sec:evaluation}

\subsection{Evaluation Questions, Setup, and Baselines}
\draftnote{Evaluate whether the methods reduce session serving time and recomputation, when benefits disappear, and which mechanisms produce the gains.
Reuse the replication setup where possible, and list all changes, including SSD model, effective read/write bandwidth, filesystem, DRAM budget, and page-cache treatment.}
\draftnote{Compare the default serving policy, CacheWise with recomputation, CacheWise with DRAM offloading, CacheWise with synchronous SSD offloading (Method 1), and each implemented optimization.
Hold CacheWise's scheduling and eviction rules fixed in the storage comparison.
Report memory and storage capacity for every configuration and make additional capacity explicit.}

\subsection{End-to-End Results}
\draftnote{Plot session serving time and useful-token goodput across increasing concurrency and memory pressure.
Also report session wall-clock duration, time to first token, tail latency, recomputed tokens, actual bytes read/written, and GPU compute time.
Include repeated runs and uncertainty estimates, and check that cache restoration preserves valid model state and output behavior.}

\subsection{Ablations and Sensitivity Analysis}
\draftnote{Add asynchronous I/O, prefetching, and cost-aware decisions individually to the synchronous SSD baseline, then evaluate the combined method.
Sweep storage bandwidth, block size, queue depth, host-buffer capacity, reuse frequency, and workload concurrency.
Include rarely reused prefixes and saturated storage to identify regressions.}

\subsection{Resource Costs and Interpretation}
\draftnote{Account for write amplification, unused prefetches, staging memory, CPU overhead, and storage occupancy alongside compute saved.
If monetary cost is reported, compare equal completed work at stated service targets and include GPU, host, and storage charges.
Connect measured outcomes to the break-even analysis in Section~\ref{sec:tradeoff}; do not infer end-to-end gains from isolated I/O timings alone.}
